\documentclass[10pt,a4paper]{amsart}
\usepackage[margin=19mm]{geometry}
\usepackage{amsmath,amssymb,mathtools}
\usepackage[scr=rsfs,cal=euler]{mathalfa}
\usepackage{xfrac}
\usepackage[unicode,hidelinks,bookmarks=false]{hyperref}
\newcommand{\ie}{i.e.\ }
\newcommand{\cf}{cf.\ }
\newcommand{\resp}{respectively\ }
\newcommand{\Subsection}{\subsection}
\providecommand{\nobreakdash}{\nobreak\hbox{-}}
\setlength{\emergencystretch}{2em}
\multlinegap0pt
\usepackage{amssymb}
\usepackage{tikz}
\usepackage[shortlabels]{enumitem}
\usepackage{xcolor}
\newtheorem{rmk}{Remark}
\newtheorem{lem}{Lemma}
\newcommand{\Cc}{\mathcal{C}}
\title{Higher-rank trees arising from polyhedral graphs}
\author{David Pask}
\date{Source: arXiv:2407.14048v3 (CC BY 4.0)}
\begin{document}
\maketitle

\noindent\emph{Source and licence.} Higher-rank trees arising from polyhedral graphs. Version arXiv:2407.14048v3 (CC BY 4.0), distributed by arXiv under the Creative Commons Attribution 4.0 International licence, http://creativecommons.org/licenses/by/4.0/, as recorded in the arXiv licence metadata for this exact version and shown on the abstract page https://arxiv.org/abs/2407.14048v3. Full licence notice: \texttt{licenses/arxiv-2407.14048-CC-BY-4.0.txt}.\par\smallskip

\noindent\textit{Editorial scope.} Counted unit: the article's lem k-colouredgraph (source0.tex lines 686--692). The article's complete original proof of it is delivered with it (source0.tex lines 694--712, one \texttt{proof} environment of 19 lines). No mathematics is written, repaired or re-derived: the statement and the proof are the article's own lines, in the article's own notation, and no retained line is substituted or re-flowed.  Retained uncounted as the source-local context the proof needs: lines 668--684; lines 670--674; lines 678--680.  Nothing was omitted from any retained span, so the package carries no omission record.  No retained line contains a citation, so the wrapper carries no bibliography and no bibliography entry.  No retained line contains a figure environment, an image-inclusion command or drawing code, so no image is delivered and the package declares no asset dependency.  Editorial formatting only, in addition: an AMS \texttt{amsart} preamble that restates the article's own declarations and macros used by the retained lines, namely the environments \texttt{rmk}, \texttt{lem} on separate counters, together with the standard packages those lines need.  The retained environments print in this document as Remark 1, Lemma 1 in order of appearance, whereas the article prints the same items with the numbering of its own full document.  The complete native source of the article as distributed in the arXiv e-print for this exact version is bundled unchanged as \texttt{sources/162560/source0.tex}.
\begin{rmk} \label{rmk:simdef}
We may expand the equivalence relation $\sim_1$ on $E^0 \sqcup E^1$, to a relation $\sim$ on all of $E^*$ as follows:  First we define a relation $\sim_2$ on $\sqcup_{i=0}^2 E^i$ as follows: 
\begin{equation} \label{eq:nested1}
x \sim_2  y \text{ if and only if } \begin{cases} x=ef, y=e'f' & e \sim_1 e', f \sim_1 f', \\
x \sim_1 y & \text{ otherwise }. \\
\end{cases}
\end{equation}

\noindent
Note that $c(e)=c(e')$ and $c(f)=c(f')$ is automatic by definition of $\sim_1$. Then for $n \ge 3$ we define a relation $\sim_n$ on $\sqcup_{i=0}^n E^i$ as follows: $x \sim_n  y$ if and only if either $x=\alpha ,y=\alpha'$ with 
\begin{equation} \label{eq:nested2}
\alpha_1\cdots \alpha_{n-1} \sim_{n-1} \alpha_1 \cdots \alpha_{n-1}',  \text{ and } \alpha_2 \cdots \alpha_n \sim_{n-1} \alpha_2' \cdots   \alpha_n', \text { or } \alpha \sim_{n-1} \alpha'  .
\end{equation}

\noindent
This leads us to a relation $\sim$ on $E^*$, whose properties we summarise below.
\end{rmk}

\begin{equation} 
x \sim_2  y \text{ if and only if } \begin{cases} x=ef, y=e'f' & e \sim_1 e', f \sim_1 f', \\
x \sim_1 y & \text{ otherwise }. \\
\end{cases}
\end{equation}

\begin{equation} 
\alpha_1\cdots \alpha_{n-1} \sim_{n-1} \alpha_1 \cdots \alpha_{n-1}',  \text{ and } \alpha_2 \cdots \alpha_n \sim_{n-1} \alpha_2' \cdots   \alpha_n', \text { or } \alpha \sim_{n-1} \alpha'  .
\end{equation}

\begin{lem}[Quotient graph] \label{lem:quotient k-colouredgraph}
Let $(E,c)$ be a $k$-coloured graph with commuting squares $\Cc$.  
Suppose that $\sim_1$ is an equivalence relation on $E^0 \sqcup E^1$
with the following properties: For $e,e' \in E^1$ if $e \sim_1 e'$ then $c(e) = c(e')$, $s(e) \sim_1 s(e')$, $r(e) \sim_1 r(e')$. Then the relation $\sim$ on $E^*$ described in Remark~\ref{rmk:simdef} above is an equivalence relation.

Set $E/{\Tiny\sim}  = (E^0/{\Tiny\sim},E^1/{\Tiny\sim},r,s)$, $c([e])=c(e)$ where $r([e])=[r(e)]$, $s([e])=[s(e)]$. Then $( E/{\Tiny\sim} , c )$ is a $k$-coloured graph with squares $\Cc / {\Tiny\sim}= \{ [e][f] {\Tiny\sim} [g][h] : ef \sim_\Cc gh \}$.
\end{lem}

\begin{proof}
That $\sim$ is an equivalence relation on $E^*$ follows by showing that each $\sim_n$, $n \ge2$ is an equivalence relation on $\sqcup_{i=0}^n E^i$. The rest is a straightforward calculation.

It follows that if $x \sim_1 y$ for $x,y \in E^0 \sqcup E^1$ then either $x,y \in E^0$ or $x,y \in E^1$ as the domain of $c$ is $E^1$, and $r(v)=s(v)=v$ if $v \in E^0$. Let $[e]_1$ denote the equivalence class containing $e \in E^1$ and suppose that $e' \in [e]_1$. Then $c(e) = c(e')$, $s(e) \sim_1 s(e')$, $r(e) \sim_1 r(e')$. Hence, the formulas
\[
c_1([e]_1) := c(e), \quad r([e]_1) = [r(e)]_1,\text{ and } s([e]_1) = [s(e)]_1
\]

\noindent
are well defined and so the quotient coloured graph $E/{\Tiny\sim_1}  = (E^0/{\Tiny\sim_1},E^1/{\Tiny\sim_1},r,s)$, $c_1([e]_1)=c(e)$ makes sense. One checks
\[
(E/{\Tiny\sim} ,c) = (E^0/{\Tiny\sim},E^1/{\Tiny\sim},r,s) ~=~ (E/{\Tiny\sim_1},c_1)   = (E^0/{\Tiny\sim_1},E^1/{\Tiny\sim_1},r,s)
\]

\noindent
where $c([e])=c(e)=c([e]_1)$ as the equivalence relations $\sim_n$, $n \ge 2$ are nested (cf. \eqref{eq:nested1}, \eqref{eq:nested2}). 

Suppose $ef \sim_\Cc gh$ in $(E,c)$, then we claim that $[e][f] \sim_\Cc [g][h]$ in $(E/{\Tiny\sim},c)$. By \eqref{eq:nested1} we have $[ef]_2 = [e]_1 [f]_1$, and similarly $[fg]_2 = [f]_1 [g]_1$. Then $[e]_1 [f]_1 = [ef]_2 = [gh]_2 = [g]_1 [h]_1$ and so $[e]_1 [f]_1 \sim_\Cc [g]_1 [h]_1$, and the result follows by the nesting property of $\sim$.
\end{proof}

\end{document}
